\scanpage{74}
% archive foliation: 73
\unsure{fonctorielle|fonctoriellement} de $X$, mais on trouverait encore
un iso fonctoriel en $X$, $\mathcal{F}_X\simeq\underline{\omega}_{X/\mathbb{C}}^{\otimes2}$.
De $=$ dans (III). Ainsi les $2$-formes sont essentielles.

2) \deleted{\unsure{G}} Le torseur $\mathcal{T}$ dont il est question dans
II (et \unsure{encore|a fortiori} dans I) dépend du choix de
$G$ (et de $\mathcal{F}$, si on veut --- c'est alors un
torseur sous $\mathcal{F}$) mais les \inserted{torseurs associés aux} différents choix
sont can. (semi)-isomorphes, de façon précise
\diagram{\begin{tikzcd}
\mathcal{T}^{G} \arrow[r,"\sim"] & \mathcal{T}^{\varphi G}\\
\underline{\omega}_X^{\otimes2} \arrow[r,"\varphi.\mathrm{id}"] &
\underline{\omega}_X^{\otimes2}
\end{tikzcd}}
\unsure{resp. aux}
i.e. $\mathcal{T}^{\varphi G}$ est ``multiple'' de $\mathcal{T}^{G}$ par $\varphi$.
Ainsi, le choix ``le plus fin'' de $\mathcal{T}$ est $\mathcal{T}^{G_0}$
(où $G_0$ est défini en (III)), alors le $\mathcal{T}^{G}$
correspondant au \inserted{$G$} \unsure{choisi|choix} de (II) est un
multiple $12$ième de $\mathcal{T}^{G_0}$.

3) On peut se \deleted{demander}\inserted{poser} la question
de caractériser $G_0(f,g)$ \uline{en une caractéristique donnée} $p>0$.
On peut prendre les conditions
a) b) c), \inserted{\unsure{vues} comme plus haut}, plus d') que $G_0(f,g)$ n'est pas id. nul
--- \uline{cela caractérise} $G_0(f,g)$ à multiplication près
par une constante $c\in\mathbb{F}_p^{*}$.
$\bar{M}$ \unsure{précisé|précisément} si
on a fixé un schéma de base absolu
\uline{intègre} $S_0$ (pour les $S$ et les $X$), l'indétermination
est dans $\Gamma(S,\mathcal{O}_S^{\unsure{\times|*}})$.
Si $S$ n'est pas intègre, il faut supposer que, sur chaque
$\unsure{k(s)|\kappa(s)}$, l'op. induit (pour des $X$ lisses
\deleted{\unsure{sur le}}\inserted{de dim. rel. $1$} sur $k(s)$)
\inserted{n'est pas identiquement nul.}

\margin{\uline{faux?} \unsure{pour un|vrai si}
\unsure{op.|W'} \unsure{diff.} d'ordre $\leq3$)}
\figure{\begin{tikzpicture}
\draw (0,0) -- (0,3.1);
\draw (0,1.1) -- (-0.6,1.1);
\draw (-0.6,1.1) circle (0.06);
\draw (0.2,0) .. controls (0.7,-0.4) and (0.8,-0.7) .. (0.2,-0.9);
\end{tikzpicture}}

